% ------------------------------------------------------------------------
% Modelo de TCC Lic Mat SPP  - versão 2
% ------------------------------------------------------------------------

% Agradecimentos a todos os colegas que compartilharam versões de TCC pelo overleaf, esta versão foi baseada em um compilado de vários modelos disponibilizados pela plataforma




\documentclass[book, 12pt,				% tamanho da fonte
	oneside,			% para impressão em frente. Outra opção é twoside (frente,costas)
	a4paper,			% tamanho do papel. 
	% % -- opções da classe abntex2 --
	% chapter=TITLE,		% títulos de capítulos convertidos em letras maiúsculas
	% -- opções do pacote babel --
	english,			% idioma adicional para hifenização
	brazil,				% o último idioma é o principal do documento
	sumario=tradicional
	]{abntex2} % carrega configurações da abnt mais recentes


% texto matemático todo em modo de displaystyle
\everymath{\displaystyle}


% --- 
% PACOTES PARA AJUSTE e CORREÇÕES EM ABNTEX2
% --- 
\input{Pasta 5 - Configurações/Pacotes}

% --- 
% Carrega informações dos dados gerais e de autoria
% --- 
\input{Pasta 0 - Dados do Trabalho/Dados Gerais}
 
% --- 
% CAPA
% --- 
\input{Pasta 1 - Capa e Folha de Rosto/capa}
 
% --- 
% FOLHA DE ROSTO
% --- 
\input{Pasta 1 - Capa e Folha de Rosto/folha-de-rosto}

% --- 
% Carrega informações de dados gerais 
% --- 
\tipotrabalho{\tipotrabalhoescrito}
\instituicao{\instituicaotrabalho}

\autor{\discente}
\titulo{\titulotrabalho}
\data{2023}
\local{\cidadededefesa}

\orientador{\proforientador}
\coorientador{\profcoorientador}

% ---
% COMPILA O INDICE
% ---
\makeindex

\input{Pasta 5 - Configurações/ComandosEspeciais}
\begin{document}

% \DoubleSpacing  espaçamento duplo (o espaçamento padrão entre linhas é 1.5)

% Seleciona o idioma do documento (conforme pacotes do babel)
\selectlanguage{brazil}

%----------------------------------------------------------
% ----------------------------------------------------------
% ELEMENTOS PRÉ-TEXTUAIS
% ----------------------------------------------------------

\pretextual

%mais informações em: https://ctan.dcc.uchile.cl/macros/latex/contrib/abntex2/doc/abntex2.pdf

% ---
% Capa
% ---
\imprimircapa

% ---
% Folha de rosto
% ---
\imprimirfolhaderosto

% ---
% Inserir folha de aprovação
% ---
\input{Pasta 2 - Elementos Pré-Textuais/folha-de-aprovacao}

% ---
% Dedicatória
% ---
\input{Pasta 2 - Elementos Pré-Textuais/dedicatoria}

% ---
% Agradecimentos
% ---
\input{Pasta 2 - Elementos Pré-Textuais/agradecimentos}

% ---
% Epígrafe
% ---
\input{Pasta 2 - Elementos Pré-Textuais/epigrafe}


% ---
% RESUMOS
% ---

% Resumo em português
\input{Pasta 2 - Elementos Pré-Textuais/resumo}

% Resumo em inglês
\input{Pasta 2 - Elementos Pré-Textuais/abstract}

% ---
% Lista de ilustrações
% ---
\pdfbookmark[0]{\listfigurename}{lof}
\listoffigures*
\cleardoublepage  % pula para próxima página
% ---

% ---
% Lista de tabelas
% ---
\pdfbookmark[0]{\listtablename}{lot}
\listoftables*
\cleardoublepage % pula para próxima página
% ---


% ---
% Sumario
% ---
\pdfbookmark[0]{\contentsname}{toc}
\tableofcontents*
\cleardoublepage % pula para próxima página

%-------------------------------------------------------
%%%%%%%%%%  Fim dos elementos pré-textuais ------------------
%----------------------------------------------------




%-----------------------------------------------------------
% ----------------------------------------------------------
% ELEMENTOS TEXTUAIS
% ----------------------------------------------------------

\textual   %comando acrescenta numeração e cabeçalho O comando \textual identifica o início dos elementos textuais. As páginas desses elementos são numeradas com algarismos arábicos no lado direito superior ou direito/esquerdo superior caso a impressão frente e verso (opção twoside) seja acionada, conforme estabelece a ABNT NBR 14724:2011. Geralmente a “Introdução” é o
% primeiro capítulo textual. A título de coerência, a macro \mainmatter, padrão do memoir, é reescrita para que tenha o mesmo comportamento que \textual. Por isso, fique livre em escolher qualquer das macros. Porém, o uso de uma delas é
% obrigatória, para que os cabeçalhos sejam montados corretamente.



%% Insere os elementos textuais
\input{Pasta 3 - Elementos Textuais/Cap1 - Introdução}
\input{Pasta 3 - Elementos Textuais/Cap2}
\input{Pasta 3 - Elementos Textuais/Cap3}
\input{Pasta 3 - Elementos Textuais/Cap4}
\input{Pasta 3 - Elementos Textuais/Cap5}
\input{Pasta 3 - Elementos Textuais/Cap6}
\input{Pasta 3 - Elementos Textuais/Cap7}
% ----------------------------------------------------------
% ----------------------------------------------------------
\phantompart




% ----------------------------------------------------------
% ELEMENTOS PÓS-TEXTUAIS
% ----------------------------------------------------------
\postextual %carrega configuração de elementos pós textuais

% ---
% Referências
% ---
\bibliography{referencias.bib}



% ---
% Apêndices
% ---
\begin{apendicesenv}

% Imprime uma página indicando o início dos apêndices
\partapendices
%\addcontentsline{toc}{section}{Seção Não Numerada}
\input{Pasta 4 - Elementos Pós-textuais/apendices}

\end{apendicesenv}







% ---
% INDICE REMISSIVO
% ---
\phantompart
\printindex



%---------------------------------------------------------------------
\end{document}